\chapter{Project compilation}
This chapter focuses on the way actual Symfony projects are compiled and deployed using Peachpied Symfony toolkit. It introduces tools that make it possible and describes the process of their development. At the end of the chapter, an overview of provided feature of template engines interoperability is revealed.

\section{Compilation setup}
Whether we want to embed our existing Symfony application within another
 .NET Core one or just compile it for enhanced security, we always have to produce it in a form of \texttt{.dll} and supply it to the .Net Core web server. In this case, we decided to choose Kestrel web server as it is both open-source and Peachpie already contains an extension that allows Kestrel to handle requests targeting \texttt{.php} files and process them with \texttt{dll} library compiled from PHP code \citep{Peachpie}. This binding is in our case done by creating two .NET projects, one for the Symfony sources and the other for the server.

\paragraph{Binding with the web server} For that, we can use an already made template called ``Web App" provided by Peachpie. Upon usage, the template produces the \texttt{Website} and the \texttt{Server} directories. The \texttt{Server} one contains simple project with a base configuration for Kestrel server and logic for \texttt{.php} request handling. The other one is expected to contain PHP code for compilation.

Apart from the configuration files \texttt{appsettings.json} and \texttt{launchSettings.json}, and project file .csproj, the server project directory by default contains only \texttt{Program.cs} file representing an entry point of the application. It consists of two classes, the \texttt{Program} class with the \texttt{static void Main(string[] args)} function and the \texttt{Startup} class with the already discussed \texttt{ConfigureServices} and \texttt{Configure} functions. Furthermore, Listings 2.1, 2.2 and 2.3 combined represent the whole default \texttt{Program.cs}. The only difference is that the \texttt{Main} method must in the Host object configuration declare the actual Symfony project name.

\paragraph{Configuring Symfony's .msbuildproj} For the majority of pure PHP web applications, the system described will already be executable, however, few further adjustments in Website's \texttt{.msbuildproj} file have to be done in order to run our application as expected.

First, we have to modify value of the \texttt{Compile} item. Website project compiles every \texttt{.php} file by default which is clearly more than we actually need. As we supply Symfony Components via package references, we can safely exclude content of the \texttt{vendor} directory which might in case of compiling an existing Symfony project occur to exist in the \texttt{Website} directory. We must not forget about the autoloading scripts that are part of the \texttt{vendor} directory as well so we have to make an exception for them in the exclusion rule.

Secondly, we need to include all the static content of the project like configuration \texttt{.yaml} files, but images and other assets as well. This is done by setting up a \texttt{Content} item.

Thirdly, package references to the Symfony Components used by the original project need to be manually set up. This is done by copying them from the project's \texttt{composer.json} file and rewriting them to the NuGet references. For this reason the \texttt{Peachpied.Symfony.Skeleton} NuGet package was created that groups together dependencies of Symfony Skeleton project with a matching version. As this thesis focuses solely on Symfony 4.2.3, this package is prepared only in this particular version. However, automatic version resolving target was implemented as a part of provided Skeleton's \texttt{.msbuildproj} file allowing their future arrangement into folders by version names.

Finally, one must not forget about the autoloading scripts and cache directory discussed in chapter 4. For now, we will settle for knowing, that all these problems will be resolved automatically by adding package reference to the \texttt{Peachpied.Symfony.ProjectTools} NuGet package. Implementation details of this set of tools will be discussed in the following section. Bare bones Symfony application's configuration for compilation will in the end look as shown in the following Listing.

\begin{center}
\lstset{language=XML,
	morekeywords={Include, Exclude, Version, CopyToOutputDirectory}}
\begin{lstlisting}[caption={Empty-Page.msbuildproj},captionpos=b]
<PropertyGroup>
	<OutputType>Library</OutputType>
	<TargetFramework>netstandard2.0</TargetFramework>
    <AssemblyName>Empty-Page</AssemblyName>
</PropertyGroup>

<ItemGroup>
	<Compile Include="**\*.php" 
			 Exclude="vendor\**; config\bootstrap.php; 
					  .\obj\**; .\bin\**" />
	<Compile Include="vendor\composer\**\*.php; vendor\autoload.php" />
    <Content Include=".\**" 
			 Exclude="{* all non configuration files *}" 
			 CopyToOutputDirectory="PreserveNewest" />
</ItemGroup>

<ItemGroup>
    <PackageReference Include="PeachPied.Symfony.Skeleton" 
    		 Version="4.2.3" />
    <PackageReference Include="Peachpied.Symfony.ProjectTools" 
			 Version="1.0.0" />
</ItemGroup>
\end{lstlisting}
\end{center}

In order to bring UX even further, all the previously mentioned steps were compressed into two new templates provided in the \texttt{Peachpied.Symfony.Templates} NuGet package. The ``Symfony Skeleton" or ``syskeleton" template creates a doublet of base Kestrel server configuration and bare bones Symfony version 4.2.3 application with all the \texttt{.php} sources provided. The second template called ``Symfony-Empty-Project" or ``syemptyproj" creates one project with the same Kestrel configuration as the previous template, but the Symfony project produced contains only the \texttt{.msbuildproj} file without any Symfony Component package references, letting user import own sources and setup required package references as required.


\section{ProjectTools package}
Just like the mentioned \texttt{ComponentTools} hooked the custom targets in order to provide necessary configuration settings during the Symfony Components compilation, \texttt{ProjectTools} does the same thing during the Symfony project compilation in order to provide the compilation with additional \texttt{.php} files needed for actual application execution. Custom targets provided by this project are named \texttt{RestoreComposerLock}, \texttt{GenerateSymfonyAutoload} and \texttt{GenerateSymfonyCache}. Because these tasks do not manipulate with project's package references, whole project can be safely delivered in a form of NuGet package. Furthermore, we can make use of discussed NuGet's feature to import special \texttt{.targets} file into out project file. This way, we prepare all the custom targets' load-up and usage in \texttt{Peachpied.Symfony.ProjectTools.targets} and let the package manager handle their embedding into compiled project's execution order, while effectively isolating user from all the background tasks.

\paragraph{RestoreComposerLock} This task is responsible for collecting all the \texttt{lockFragments} from NuGet packages of Symfony Components and joining them together into \texttt{composer.lock} file. The data from packages are in fact extended by the data taken from project's \texttt{composer.json} file in order to provide a file that matches the one originally produced by Composer. This also required to rewrite implementation of content hashing function used in composer into C\# in order to provide value for the \texttt{content-hash} key.

This task's execution is bound by appropriate \texttt{Inputs} and \texttt{Outputs}, ensuring that it runs only when there is no \texttt{composer.lock} file present or its content is outdated compared to \texttt{composer.json}.

\paragraph{GenerateSymfonyAutoload} Things get slightly more complicated with this task due to requirements to use composer itself to produce these file for us. The implementation we settled for includes the whole \texttt{composer.phar} file as a static asset of the NuGet package and calls it on task's execution with the \texttt{php composer.phar dump-autoload} command.

Unfortunately, quite severe problem emerged regarding composer's \texttt{dump-autoload} requirements. For its proper execution, this task needs the \texttt{installed.json} file to be located under the \texttt{vendor/composer} path in the project directory. Thankfully, this file is nothing more than a concatenation of each package's \texttt{lockFragment} and thus can be easily created using previously generated \texttt{composer.lock}. On the other hand, it does also require for all the \texttt{.php} source codes of the Symfony Components that are autoloaded to be physically present under the respective subpath of the \texttt{vendor} directory. We have this resolved in the Symfony Components compilation as well and we can make Components' NuGet packages to copy these files for us by setting the \texttt{RestoreVendor} property to \texttt{true}. This is done once again by default inside of the \texttt{Peachpied.Symfony.ProjectTools.targets} resulting in the background recovery of the \texttt{vendor} directory in case it is missing. We can see now that this approach introduces serious time overhead during the first and each subsequent project compilation that results in \texttt{vendor} folder recovery. For maximal efficiency, there are constraints ensuring this recovery gets triggered only for Symfony Components whose subdirectory of \texttt{vendor} directory is missing. Were we to successfully compile composer into .NET, we would be able to omit the whole recovery phase and supply compiled \texttt{.php} sources directly from NuGet packages.

\paragraph{RestoreSymfonyCache} After autoloading scripts are produced, this task undirectly uses them with \texttt{src/Kernel.php} to produce Symfony cache. This is usually done by calling the \texttt{boot} method of the \texttt{Kernel} class instance. Additionally, routing cache can be produced with the \texttt{getMatcher} method of the \texttt{Router} class. The process is complicated due to fact that these classes load configuration files stored under the \texttt{config} directory which are not only \texttt{.yaml} files but there are some configuration \texttt{.php} files as well. Because of their dynamic content, it is not possible to precompile them, and thus the only possible way to handle them from .NET would be to write a custom parser of these \texttt{.php} files. However, as we already depend on PHP being installed on the user's machine, nothing stops us from using these methods natively from PHP. For this reason, a simple PHP script handling cache generating was written, included with the \texttt{composer.phar} in the \texttt{ProjectTools} NuGet package and is triggered during task's execution.

Furthermore, the constructor of the \texttt{Kernel} class accepts parameters identifying the current environment and thus influencing produced cache. These parameters are \texttt{string} value for \texttt{\$environment} and \texttt{boolean} value for \texttt{\$debug}. Values for these parameters are passed to the PHP script from the task where they are obtained by manual parsing either from the \texttt{.env} file or the \texttt{.appsettings.json} file. As is discussed in greater detail later on, Symfony uses the \texttt{.env} and \texttt{.env.[environment]} files for environmental variables configuration, where .NET Core uses the \texttt{appsettings.json} and \texttt{appsettings.[environment].json}.
 
 
\section{Peachpied.Symfony.AspNetCore}
With the configuration as shown in the Listing 6.1, the application will successfully compile into a \texttt{.dll}. \texttt{Program.cs} file in the Server project provided by one of Symfony templates will then start a web server and upon accessing the application on default port, content of the \texttt{public/index.php} file will be displayed. In order for everything to work as expected, the Server project uses \texttt{Peachpied.Symfony.AspNetCore} NuGet package and its \texttt{UseSymfony} method. Moreover, many additional efforts had to be made in misbehaviour errors identification, and by the PeachPie development team in particular, during Peachpie-related problems resolving. Significant errors encountered will be discussed later.

This package is intended as a helper class for the Server project, encapsulating the Symfony specific server configuration and implementation of additional interoperability features. Both are provided via the Application programming interface (API), following the standard ASP.NET Core \texttt{Use[midlewareName]} naming conventions. The public API is provided by the \texttt{RequestDelegateExtension} class and provides the \texttt{UseSymfony} method discussed in the following paragraph.

\paragraph{UseSymfony method} As the method's name suggests, it adds middleware to the pipeline that handles requests targeting Symfony application. Everything that needs to be provided is a relative path to the root of the project.

Aside from calling the \texttt{UsePhp} method provided by the \texttt{Peachpie.AspNetCore.Web} which implements the actual middleware, the method configures static file serving middleware, further discussed URL rewriting and environmental variables load-up, and optionally overrides Symfony's \texttt{bootstrap.php} which is responsible for environmental variables loading in Symfony.

The content of the \texttt{bootstrap.php} file can be found in the Listing 6.2 without original comments that have been omitted for the sake of increased readability. The fragment of the code performing the overriding is captured in the Listing 6.3. This is done by accessing the \texttt{Context} object's script table that performs mapping from \texttt{.php} files to compiled .NET classes, and overwriting path to class stored for the \texttt{bootstrap.php} file. It is also worth of notice, that our \texttt{BootstrapMain} class effectively substitutes original file's additional functionality and overrides just the environmental configuration.

\begin{center}
\lstset{language=PHP,
	morekeywords={new, null, throw, class, use}}
\begin{lstlisting}[caption={bootstrap.php},captionpos=b]
use Symfony\Component\Dotenv\Dotenv;

require dirname(__DIR__).'/vendor/autoload.php';

if (is_array($env = @include dirname(__DIR__).'/.env.local.php')) {
	$_SERVER += $env;
	$_ENV += $env;
} elseif (!class_exists(Dotenv::class)) {
	{* Runtime exception *}};
} else {
	(new Dotenv())->loadEnv(dirname(__DIR__).'/.env');
}

$_SERVER['APP_ENV'] = $_ENV['APP_ENV'] = ($_SERVER['APP_ENV']
	?? $_ENV['APP_ENV'] ?? null) ?: 'dev';
$_SERVER['APP_DEBUG'] = $_SERVER['APP_DEBUG']
	?? $_ENV['APP_DEBUG'] ?? 'prod' !== $_SERVER['APP_ENV'];
$_SERVER['APP_DEBUG'] = $ENV['APP_DEBUG'] = (int) $_SERVER['APP_DEBUG']
	|| filter_var($_SERVER['APP_DEBUG'], FILTER_VALIDATE_BOOLEAN)
		? '1' : '0';
\end{lstlisting}
\end{center}	

\begin{center}
\lstset{language=[Sharp]C,
	morekeywords={var}}
\begin{lstlisting}[caption={Context file mapping overriding},captionpos=b]
static SymfonyConfig bootstrapConfig;

static PhpValue BootstrapMain(
	Context ctx,
	PhpArray locals,
	object @this,
	RuntimeTypeHandle self
) {
	ctx.Include("vendor", "autoload.php", true);
	Apply(ctx, bootstrapConfig);
	return 0;
}
...
// bootstrap.php env loading overriding
string bootstrapPath = "config\\bootstrap.php";
Context.MainDelegate md = newContext.MainDelegate(BootstrapMain);
Context.DeclareScript(bootstrapPath, md);
\end{lstlisting}
\end{center}


\paragraph{Environmental variables} Because Peachpied Symfony performs project migration to the ASP.NET Core environment using the .NET based web server, it can be expected that users will want to set up environmental variables in a .NET-specific way with the use of previously mentioned \texttt{appsettings.json} configuration file. On the other hand, this way of configuring should not be compulsory as there might also be users who would prefer to simply compile existing Symfony project without the need to worry about environmental variables letting the application use Symfony's configuration already present in the code. This functionality is delivered by the \texttt{SymfonyConfig} and the \texttt{SfConfigurationLoader} classes. The first one represents a definition for environmental variables and their default values. Therefore, it consists solely of properties. The second class contains methods for configuration values parsing, default configuration loading and \texttt{appsettings.json} files processing.


\paragraph{Runtime problems} Probably the most significant problem encountered was due to the differences in server root directory between Symfony and .NET. In Symfony, it is expected that a server operates from the root directory set to \texttt{public/}, therefore, all paths to assets like styles, scripts and images are set relative to this folder. In our .NET application however, we set the server root to the parent directory of the Symfony project. We have to do it this way in order to be able to reach assets from both the applications. Peachpie's \texttt{UsePhp} midlleware simplifies this problem partially by locating Symfony projects' sources for us, however, it does not help us with the problem of assets that are expected to be relative to the \texttt{public} directory.

This problem was resolved by adding URL rewriting into the \texttt{UseSymfony} method. URL rewriting is a technique used to modify request URLs based on one or more predefined rules \citep{microsoft.docs}. It creates an abstraction between resource locations and their addresses so that the locations and addresses are not tightly linked. \citep{microsoft.docs}. In ASP.NET Core, the \texttt{UseRewriter} middleware serves this cause and for that accepts instance of the \texttt{RewriteOptions} class as an argument. The actual implementation of the URL rewriting in Symfony can be found in the following Listing.

\begin{center}
\lstset{language=[Sharp]C,
	morekeywords={var}}
\begin{lstlisting}[caption={Symfony URL rewritting middleware},captionpos=b]
var options = new RewriteOptions()
    .AddRewrite(
        @"((.*\/)*([^\/\?&]*\.[^\/\?&]*))",
        bootstrapConfig.PublicDir + "/$1",
        skipRemainingRules: true
    )
    .AddRewrite(
        @"^(?!"+ bootstrapConfig.PublicDir + "/)(.*)$", 
        bootstrapConfig.PublicDir, 
        skipRemainingRules: true
    );

app.UseRewriter(options);
\end{lstlisting}
\end{center}

Regarding the Peachpie-related errors we will mention invalid handling of the \texttt{new ReflectionProperty("Exception", "trace")} expression. This was done due to lack of formal documentation of private \texttt{\$trace} property of the \texttt{Exception} class. For the moment, this enabled the application to run only with the \texttt{APP\_DEBUG} environmental variable set to \texttt{false}. However, the error was later on resolved by an update of PeachPie compiler and therefore the debug mode is fully supported at the moment. Along some less complicated errors we can name incorrect \texttt{realpath} function return value implementation, where the function returned \texttt{null} instead of \texttt{false} on failure.


\section{Additional interoperability}
In addition to the already mentioned \texttt{UseSymfony} method, the \texttt{Peachpied.Symfony.AspNetCore} API provides two classes that display some of the interoperability possibilities brought by the Symfony compilation. These classes are named \texttt{TwigRenderer} and \texttt{RenderRazorService}, and both are stored in the \texttt{Templating} namespace of the \texttt{Peachpied.Symfony.AspNetCore} NuGet package.

\paragraph{Twig into Razor} The \texttt{TwigRenderer} class provides very promising tools to embed Twig templates within Razor templates. Furthermore, it can be all done without the \texttt{Peachpied.Symfony.Skeleton} NuGet package, but only through using two additional packages, the \texttt{Twig.Twig} package and its single dependency, the \texttt{Symfony.Polyfill-mbstring} package.

In order to make this process as simple as possible, the \texttt{PeachPied.Symfony.AspNetCore.Templating} namespace provides the \texttt{TwigRenderer} class, with the static \texttt{RenderTwig} function that can be directly accessed and used from a Razor template. Furthermore, additional \texttt{DataToPhp} function was implemented in order to provide seamless means for converting a Razor data into the Twig comprehensible form. Finally, the \texttt{MergePhpArrays} function offers merging of converted data into a single unit that can be supplied to the \texttt{RenderTwig} function afterwards.

The result of the Twig rendering process is stored as a \texttt{string} and returned into a Razor template. After that, it can be further processed or directly embedded into the code using the \texttt{@Html.Raw()} method.


\paragraph{Razor into Twig} Due to a certain level of complexity Razor into a \texttt{string} rendering brings, the whole functionality is separated into three functional units, each placed into its own class. 

The first unit, the \texttt{RenderRazorService} class, provides a functionality for Razor into a \texttt{string} rendering, accessible via the \texttt{RenderToString} function in the \texttt{PeachPied.Symfony.AspNetCore.Templating} namespace. However, such feature requires access to the objects like the \texttt{RazorViewEngine} or the \texttt{ServiceProvider}. To take the advantage of the Dependency Injection of Asp.Net Core, the class needs to be registered as a service in the server's \texttt{Startup} class, within the \texttt{ConfigureServices} method, implementing the \texttt{IRenderRazorService} interface provided. This way, the application provides the necessary objects for us, creating an instance of the \texttt{RenderRazorService} that can further accessed through the \texttt{ServiceProvider}.

The second functional unit is called the \texttt{RSEnvironment}. This class forms a wrapper around the Twig's \texttt{Environment} class. The wrapper is located inside the \texttt{Twig} namespace and thus directly accessible from a Symfony project. Its purpose is to instantiate the standard Symfony \texttt{Environment} class and extend it with an additional function for Razor templates rendering. This is done by using the Twig \texttt{Environment}'s \texttt{addFunction} extension method. This way, a new \texttt{render\_razor} function is defined within this instance of the \texttt{Environment} that is bind to an actual Razor rendering function situated in the third functional unit. For this reason, the wrapper is named the Razor Supporting Environment, \texttt{RSEnvironment} in short.  

The last functional unit is titled the \texttt{RazorRendererBridge}, as it provides bridging between the \texttt{RSEnvironment} and the \texttt{RazorRendererService}. Each call of the \texttt{render\_razor} function within a Twig template results in the \texttt{RenderRazor- Bridge} class instantiating with the current PHP \texttt{Context}, and the \texttt{renderRazor} function invocation. This method converts the data provided into the Razor comprehensible formate, extracts the current \texttt{HttpContext} from the Peachpie's PHP \texttt{Context} provided on the instantiation, uses it to access the \texttt{RazorRenderService} and returns output from its \texttt{RenderToString} method, invoked with the data provided. 

Again, the output of Razor rendering is stored as a \texttt{string} and returned into Twig. After that, it can be further processed or directly embedded using the \texttt{|raw} filter.
